\chapter{What an observer can recover}\label{ch:access}

The previous chapters studied the topology of the history complex $K_A$ directly. An observer of a generative system rarely sees $K_A$; she sees a count of histories, the degrees of the swap graph, the sizes of neighbourhoods. An \emph{observation} is a function of the admissible set $A$ (equivalently of the regime $(\lambda,m)$), and recovery questions are questions about which functions of $A$ are functions of which others. This chapter fixes the observation ladder, proves the exact relations among its levels, develops the accessibility profile of the swap graph, and proves the radius-two identity that governs everything an observer learns from two-step neighbourhoods. Whether these observations determine \emph{topology} is the subject of \cref{ch:separation}.

\section{The observation ladder}\label{sec:ladder}
Throughout, $k=6$ unless stated, $A\subseteq S_k$ is an admissible set, $G_A$ its legal-swap graph (\cref{def:swap}) and $K_A$ the permutohedral complex (\cref{def:KA}). Write $d(w)$ for the degree of $w$ in $G_A$, $B_r(w)$ for the graph ball of radius $r$ and $N=|A|$.
\begin{definition}[observation levels]\label[definition]{def:ladder}\leavevmode
\begin{itemize}[nosep]
\item $O_N$: the history count $N=|A|$.
\item $O_h$: $N$ together with the radius-one and radius-two ball-size histograms $h_1=\{d(w)\}_w$ and $h_2=\{|B_2(w)|\}_w$ (multisets over $w\in A$).
\item $O_G$: the isomorphism class of the uncoloured graph $G_A$.
\item $O_C$: the isomorphism class of the swap-coloured graph (edge $w\sim ws_i$ coloured $i$, up to permutation of colours that is a diagram automorphism of the Coxeter graph, i.e.\ identity or reversal).
\item $O_A$: the labelled arrangement set $A$ itself (a subset of $S_k$).
\end{itemize}
A \emph{target} $Y$ is a function of $A$ (for instance $Y=\mathbf 1[\beta_2(K_A;\F)>0]$). We say $O$ \emph{recovers} $Y$ on a domain $\mathcal D$ of sets if $Y$ is constant on every fibre of $O$ restricted to $\mathcal D$.
\end{definition}

\begin{figure}[ht]\centering
\begin{tikzpicture}[font=\small,
 okn/.style={minimum width=1.9cm,minimum height=0.75cm,draw,fill=accent!16,align=center},
 badn/.style={minimum width=1.9cm,minimum height=0.75cm,draw,fill=red!12,align=center},
 unkn/.style={minimum width=1.9cm,minimum height=0.75cm,draw,fill=gray!10,align=center}]
 \foreach \i/\l in {0/{$O_N$},1/{$O_h$},2/{$O_G$},3/{$O_C$},4/{$O_A$}}{\node[font=\bfseries] at (2.0*\i,1.5) {\l};}
 \node[anchor=east] at (-1.1,0.8) {$\beta_0$};
 \node[anchor=east] at (-1.1,0.0) {$\beta_1$};
 \node[anchor=east] at (-1.1,-0.8) {$\beta_2>0$};
 \node[anchor=east] at (-1.1,-1.6) {anchored state};
 \node[unkn] at (0,0.8) {not claimed};\node[unkn] at (2,0.8) {not claimed};
 \foreach \i in {2,3,4}{\node[okn] at (2.0*\i,0.8) {determines};}
 \node[unkn] at (0,0.0) {not claimed};\node[unkn] at (2,0.0) {not claimed};
 \foreach \i in {2,3,4}{\node[okn] at (2.0*\i,0.0) {determines$^\ast$};}
 \foreach \i in {0,1}{\node[badn] at (2.0*\i,-0.8) {fails};}
\node[unkn] at (4.0,-0.8) {open};
 \foreach \i in {3,4}{\node[okn] at (2.0*\i,-0.8) {determines};}
 \foreach \i in {0,1,2}{\node[badn] at (2.0*\i,-1.6) {fails};}
\node[unkn] at (6.0,-1.6) {open};
 \node[okn] at (8,-1.6) {determines};
\end{tikzpicture}
\caption{The observation ladder against four targets. Blue: the level determines the target on all admissible sets; red: distinct admissible sets with equal observation and different target exist; grey (open): neither is claimed. $^\ast$For prefix-defined $A$ and $k\le8$ (\cref{thm:skeleton}). For $\beta_2>0$ the uncoloured graph fixes everything but $\operatorname{rank}\partial_3$ (\cref{thm:third-boundary}); whether it determines $\beta_2$ on realizable systems is not settled here. For the anchored state, failure at $O_G$ (\cref{prop:uncoloured-fails}) implies failure at the coarser levels. Whether radius-two data $O_h$ still recover $\beta_2>0$ on the arithmetic census is a different question, answered in \cref{ch:separation}.}\label{fig:ladder}
\end{figure}

\begin{theorem}[the observation ladder; analytic]\label[theorem]{thm:obs-ladder}
Each level is a function of the next: $O_A\Rightarrow O_C\Rightarrow O_G\Rightarrow O_h\Rightarrow O_N$ (that is, $O_N$ is a function of $O_h$, which is a function of $O_G$, and so on). Consequently, for every target $Y$ and every probability measure on a finite domain,
\[
H(Y\mid O_A)\le H(Y\mid O_C)\le H(Y\mid O_G)\le H(Y\mid O_h)\le H(Y\mid O_N)\le H(Y).
\]
If $Y=\mathbf 1[\beta_j(K_A;\F)>0]$ then $H(Y\mid O_A)=0$ since $Y$ is a function of $A$, and $H(Y\mid O_C)=0$ for every $A\subseteq S_k$ (\cref{thm:swap-reconstruction}: the coloured graph determines $K_A$ up to isomorphism). Moreover $\beta_0$ is a function of $O_G$ always, and $\beta_1$ is a function of $O_G$ for prefix-defined $A$ when $k\le8$ (\cref{thm:skeleton}); for $j=2$ the uncoloured graph determines $\beta_2$ only up to $\operatorname{rank}\partial_3$ (\cref{thm:third-boundary}), and whether graph-isomorphic realisable sets can differ in $\beta_2$ is not settled here.

\end{theorem}
\begin{proof}
The ball-size histograms and $N$ are isomorphism invariants of the graph, so $O_h$ is a function of $O_G$; $O_G$ is a function of $O_C$ (forget the colours) and $O_C$ of $O_A$. The entropy inequalities are $H(Y\mid f(O))\ge H(Y\mid O)$ for any function $f$. The remaining statements are \cref{thm:swap-reconstruction}, the fact that components of $G_A$ are the connected components of $K_A$, \cref{thm:skeleton} (the $2$-skeleton, hence $f_0,f_1,f_2,\rho_1,\rho_2$ and $\beta_1=f_1-\rho_1-\rho_2$, is determined by the uncoloured graph under the cube-product property for $k\le8$), and \cref{thm:third-boundary}.
\end{proof}
\subsection*{The anchored state is not a function of uncoloured data}
For the two anchor families (\cref{ch:topology}) define the \emph{anchored hole state} $Y_{\mathrm{anc}}(m)\in\{1,0,\mathrm{NA}\}$: $1$ if the generator of $H_2(K_A(m_0))$ maps to a nonzero class in $H_2(K_A(m))$, $0$ if it maps to zero, and $\mathrm{NA}$ if $A(m)\not\supseteq A(m_0)$ so that the map is undefined.
\begin{proposition}[uncoloured data can fail to recover the anchored state; finite-exact]\label[proposition]{prop:uncoloured-fails}
For $\lambda_A=(2,3,4,5,8,9)$: the sources $m=2160$ (state $\mathrm{NA}$) and $m=2880$ (state $0$) have isomorphic uncoloured swap graphs, equal ball histograms, equal Betti numbers and equal $f$-vectors, and are the same abstract complex under the rank relabelling $\sigma=(2,0,1,3,5,4)$ of the admissible families; only the labelled arithmetic data ($O_A$) separate them. For $\lambda_B$ no such pair exists.

\end{proposition}
The point is conceptual: the anchored state depends on the identification of a specific cycle across sources, which is labelled information, whereas $\mathbf 1[\beta_2>0]$ depends only on the complex.

\section{The accessibility profile}\label{sec:profile}
Define the \emph{recombination information at radius $r$} for the uniform measure on $A$:
\[
I_r=\frac1N\sum_{w\in A}\log_2|B_r(w)|=\frac1N\log_2P_r,\qquad P_r=\prod_{w\in A}|B_r(w)|\in\mathbb N .
\]
$I_r$ is the mean of the Hartley entropy of the number of histories reachable from $w$ in at most $r$ swaps.
\begin{theorem}[accessibility profile; analytic]\label[theorem]{thm:profile}
Let $G$ be the swap graph of a nonempty finite fibre with $N$ vertices, $\beta_0$ components, $E$ edges and let $\rho$ be the largest eccentricity.
\begin{enumerate}[label=(\alph*)]
\item $I_0=0$ and $I_r\le I_{r+1}$.
\item $I_r<I_{r+1}$ for every $r<\rho$, and $I_r=I_\rho$ for all $r\ge\rho$.
\item $I_r=H(W\mid C,T=t)$ iff $r\ge\rho$ (uniform model), and always $I_r\le H(W\mid C,T=t)$.
\item (Additivity.) If the fibre is a disjoint union of pieces $p$ of sizes $n_p$ then $I_r=\sum_p(n_p/N)I_r(p)$, and Betti numbers add.
\item (First-step law.) $|B_1(w)|=1+d(w)$ and $I_1\le\log_2(1+2E/N)=\log_2\bigl(3+2(\beta_1(G)-\beta_0)/N\bigr)$, equality iff $G$ is regular. Hence $I_1\ge\log_23$ implies $\beta_1(G)\ge1$ (a forest has $I_1<\log_23$), and quantitatively $\beta_1(G)\ge\beta_0+N(2^{I_1}-3)/2$.
\item $\beta_1(G)$ here is the \emph{graph} cycle rank; the complex satisfies $\beta_1(K_A)\le\beta_1(G)$ (two-cells only kill cycles).
\end{enumerate}

\end{theorem}
\begin{proof}
(a) $B_0(w)=\{w\}$ and $B_r\subseteq B_{r+1}$. (b) $B_{r+1}(w)=B_r(w)$ iff no vertex of the component of $w$ lies at distance exactly $r+1$, iff $r\ge\operatorname{ecc}(w)$; so $I_{r+1}-I_r=\frac1N\sum_w\log_2(|B_{r+1}(w)|/|B_r(w)|)>0$ iff some $w$ has $\operatorname{ecc}(w)>r$, i.e.\ $r<\rho$. (c) For $r\ge\rho$ balls are components: $I_r=\sum_i(n_i/N)\log_2n_i=H(W\mid C,T=t)$ (\cref{prop:class-entropy}); for $r<\rho$ some ball is a proper subset of its component and the sum is strictly smaller. (d) Balls are computed inside a piece, so the sum splits; Betti numbers of a disjoint union of complexes add. (e) The graph is simple; Jensen's inequality for the concave $\log_2$ gives $\frac1N\sum\log_2(1+d(w))\le\log_2\frac1N\sum(1+d(w))=\log_2(1+2E/N)$, with equality iff all $1+d(w)$ agree. The cycle space has dimension $E-N+\beta_0$, so $E=N-\beta_0+\beta_1(G)$ and $1+2E/N=3+2(\beta_1(G)-\beta_0)/N$. For a forest this is $3-2\beta_0/N<3$. (f) Two-cells can only bound 1-cycles.
\end{proof}
\begin{remark}
Statement (e) says that sufficiently large radius-one information \emph{certifies some graph cycle}. It does not certify a hole in the complex: graph cycles are filled by squares and hexagons, and whether a cycle is filled is the content of \cref{ch:topology,ch:atlas}.
\end{remark}

\section{The radius-two identity}\label{sec:radius2}
Define for $w\in A$ the quantities $s(w)=\sum_{u\sim w}(d(u)-1)$ (the number of non-backtracking two-step walks starting at $w$) and $q(w)$, the number of square faces of $K_A$ through $w$. Let $E$ be the number of edges and $Q=\tfrac14\sum_wq(w)$ the number of square faces.
\begin{theorem}[radius-two identity; analytic]\label[theorem]{thm:radius2}
For every $w\in A$ and every $A\subseteq S_k$,
\[
|B_2(w)|=1+d(w)+s(w)-q(w).
\]

\end{theorem}
\begin{proof}
\emph{Step 1: bipartiteness.} Each adjacent-position swap changes the sign of a permutation, so the permutohedron graph, and hence $G_A$, is bipartite. Every non-backtracking walk $w$--$u$--$v$ with $v\ne w$ therefore ends at a vertex at distance exactly $2$ (distance $1$ would give a triangle), and every vertex at distance $2$ arises. The number of such walks is $s(w)$.
\emph{Step 2: common neighbours in the full permutohedron.} Let $v=ws_as_b$ with $a\ne b$ (positions of the two swaps). If $|a-b|\ge2$ the swaps commute and $v=ws_bs_a$ as well, so $v$ has exactly the two middle vertices $ws_a,ws_b$ in $P_k$; if $|a-b|=1$, $s_as_b$ is a $3$-cycle whose only factorisation as a product of two adjacent transpositions is $(a,b)$, and $v$ has exactly one middle vertex. (A common neighbour of $w$ and $v$ is $ws_c$ with $s_cs_d=s_as_b$ for some $d$.)
\emph{Step 3: restrict to $A$.} A distance-$2$ vertex $v\in A$ is reached by $c(v)\in\{1,2\}$ walks inside $G_A$, and by two iff $v=ws_as_b$ with $|a-b|\ge2$ and \emph{both} middle vertices lie in $A$, i.e.\ iff $w,ws_a,v,ws_b$ are the four vertices of a commuting square of $P_k$ contained in $A$, i.e.\ a square face of $K_A$ through $w$. Distinct squares through $w$ have distinct opposite vertices (since $c\le2$). So the number of vertices at distance $2$ is $\sum_vc(v)-\#\{v:c(v)=2\}=s(w)-q(w)$. Adding $w$ and its $d(w)$ neighbours gives the formula.
\end{proof}
\begin{corollary}[moment identities; analytic]\label[corollary]{cor:moments}
Write $S_{d2}=\sum d^2$, $S_{s2}=\sum s^2$, $S_{q2}=\sum q^2$, $S_{ds}=\sum ds$, $S_{dq}=\sum dq$, $S_{sq}=\sum sq$, $S_{B}=\sum|B_2|$, $S_{BB}=\sum|B_2|^2$ and $S_s=\sum s=S_{d2}-2E$.
\begin{enumerate}[label=(\roman*)]
\item $S_B=N+S_{d2}-4Q$; equivalently $Q=(N+S_{d2}-S_B)/4$.
\item For two sets with equal $N$ and equal radius-one histogram (hence equal $E$ and $S_{d2}$): $S_B$ equal $\iff$ $Q$ equal. The mean of $|B_2|$ carries no information beyond $(N,h_1)$ and the square count.
\item $S_{BB}=N+S_{d2}+S_{s2}+S_{q2}+4E+2S_s-8Q+2S_{ds}-2S_{dq}-2S_{sq}$.
\item Anything beyond first order in $|B_2|$ enters only through the second-order walk sums $S_{s2}$, $S_{ds}$, $S_{sq}$, $S_{dq}$, $S_{q2}$: sums over pairs of two-step walks from a common start.
\end{enumerate}

\end{corollary}
\begin{proof}
Sum \cref{thm:radius2} over $w$: $\sum d=2E$, $\sum s=\sum_u d(u)(d(u)-1)=S_{d2}-2E$ (each $u$ is counted once for each of its $d(u)$ neighbours, each term $d(u)-1$), $\sum q=4Q$; so $S_B=N+2E+S_{d2}-2E-4Q$, giving (i); (ii) follows. Squaring the identity and summing gives (iii): $(1+d+s-q)^2=1+d^2+s^2+q^2+2d+2s-2q+2ds-2dq-2sq$. (iv) is a reading of (iii).
\end{proof}
\begin{remark}[what the identity does and does not say]
The identity shows that fine radius-two data determine $N$, $E$ and the number $Q$ of commuting squares, and contain no \emph{direct} term for hexagonal (braid) faces, whose boundary $6$-cycles first affect balls of radius three. This is an interpretation: $s(w)$ and $q(w)$ are graph quantities whose joint distribution could correlate with hexagons. The witnesses of \cref{prop:witnesses} show that any such correlation does not recover topology in general, while on the arithmetic census radius-two data do separate the hole indicator (\cref{obs:residual}); why is open.
\end{remark}
